% =============================================================================
% DISCUSSION - CITA Paper in ALKALI Format
% =============================================================================

\vspace{-2mm}
\section{Discussion}
\label{sec:discussion}

We analyze CITA's strengths and limitations, discuss ethical considerations, and outline future research directions.

\vspace{-2mm}
\subsection{Strengths and Limitations}
\label{sec:strengths_limitations}

\begin{table}[h!]
\centering
\small
\begin{tabular}{@{}p{3.8cm}p{3.8cm}@{}}
\toprule
\textbf{Strengths} & \textbf{Limitations} \\
\midrule
Dynamic behavioral switching via natural language & Single base model (Llama-3.1-8B) \\
No reward model required & English-only evaluation \\
Mandatory KL prevents collapse & 10 instruction types tested \\
Outperforms DPO in 3/5 metrics & Requires instruction at inference \\
Publicly available ISD benchmark & Longer context windows needed \\
\bottomrule
\end{tabular}
\caption{Summary of CITA's strengths and current limitations.}
\label{tab:strengths_limitations}
\vspace{-4mm}
\end{table}

\paragraph{Model Generalization.} We evaluate only Llama-3.1-8B. Future work should validate CITA on diverse architectures (GPT, Mistral, Gemma) and scales (7B--70B) to establish generality.

\paragraph{Instruction Coverage.} The 10 instruction types in ISD cover common alignment dimensions but may not capture all real-world policy requirements. Extending to domain-specific instructions (medical, legal, educational) is important.

\paragraph{Multilingual Alignment.} All experiments are English-only. Instruction-conditioned alignment in multilingual settings raises questions about cross-lingual instruction transfer.

\paragraph{Inference Overhead.} CITA requires alignment instructions in every prompt, increasing context length by 10--40 tokens. This is minimal overhead for modern LLMs but may matter for latency-critical applications.

\vspace{-2mm}
\subsection{Why Mandatory KL?}
\label{sec:why_kl}

A natural question is: \textit{Why does CITA require mandatory KL while DPO works with implicit regularization?}

\paragraph{Hypothesis.} When conditioning on instructions, the model must maintain multiple behavioral modes simultaneously. Without explicit KL:
\begin{enumerate}[leftmargin=*,itemsep=1pt]
    \item The model may collapse to a single dominant mode (ignoring instructions)
    \item Preference margins can grow unbounded, destabilizing training
    \item The model may overfit to specific instruction-prompt combinations
\end{enumerate}

The mandatory KL term acts as a ``behavioral anchor,'' keeping the model close enough to the reference policy to respond appropriately to different instructions rather than memorizing fixed responses.

\paragraph{Empirical Evidence.} In ablation studies (Appendix~\ref{sec:appendix_ablations}), removing KL or reducing $\lambda$ below 0.0001 caused:
\begin{itemize}[leftmargin=*,itemsep=1pt]
    \item Reward margin collapse (from 7.5 to $<$4.0)
    \item Reduced instruction sensitivity (ISD scores dropped 20--30\%)
    \item Training instability (loss spikes, gradient explosions)
\end{itemize}

\vspace{-2mm}
\subsection{Ethical Considerations}
\label{sec:ethics}

\paragraph{Dual-Use Concerns.} Instruction-conditioned alignment could be misused to make models \textit{less} safe by providing adversarial instructions. Mitigations include:
\begin{itemize}[leftmargin=*,itemsep=1pt]
    \item Hierarchical instruction handling (system instructions override user instructions)
    \item Instruction validation/filtering at deployment
    \item Constitutional constraints~\cite{bai2022constitutional} as non-negotiable baselines
\end{itemize}

\paragraph{Alignment Washing.} Organizations might claim ``aligned'' behavior while using permissive instructions. Transparency about instruction policies and third-party auditing are essential.

\paragraph{User Autonomy vs. Safety.} CITA enables user-specified alignment, raising questions about who controls the alignment policy. We advocate for tiered systems where users can adjust within bounds set by deployers, who operate within regulatory constraints.

\vspace{-2mm}
\subsection{Future Research Directions}
\label{sec:future_work}

\paragraph{Hierarchical Instructions.} Multi-level instruction handling where system-level policies override user preferences, useful for enterprise deployments with governance requirements.

\paragraph{Instruction Compositionality.} Training models to combine multiple instructions (e.g., ``be concise AND professional'') in semantically meaningful ways.

\paragraph{Instruction Robustness.} Evaluating resistance to instruction injection attacks and developing defenses for instruction-conditioned models.

\paragraph{Theoretical Foundations.} Understanding the learning dynamics of instruction-conditioned preference optimization, including convergence guarantees and sample complexity.

\paragraph{Multi-Modal CITA.} Extending instruction-conditioned alignment to vision-language models where behavioral instructions may span modalities.
